\hypertarget{chapter-7-conditionals-loops-and-key-value-interfaces}{%
\chapter{Conditionals, Loops, and Key-Value Interfaces}\label{chapter-7-conditionals-loops-and-key-value-interfaces}}

A macro that always does the same thing with a fixed set of arguments is only useful up to a point. Most of the macros a superuser ends up writing for a large project need to behave differently depending on what they are called with, or need to process a variable amount of input, or need to expose a set of independent options an author can turn on or off without the macro's calling convention becoming an unreadable string of positional arguments. This chapter covers the tools for all three: conditionals, loops, and key-value interfaces, the last of which is the pattern most worth learning well, since it is what makes a package's public interface usable by someone other than its author.

\hypertarget{conditionals}{%
\section{Conditionals}\label{conditionals}}

TeX's native conditionals, \texttt{\textbackslash{}ifnum}, \texttt{\textbackslash{}ifdim}, \texttt{\textbackslash{}ifx}, and the small family of related primitives, are expandable, in the sense introduced in the previous chapter: they can appear inside a context requiring full expansion, and they produce one branch's tokens or the other's rather than executing a statement. \texttt{\textbackslash{}ifx} in particular is worth understanding well beyond its common use for comparing two macros for equality, since it compares the current meaning of two tokens, which makes it the standard tool for testing whether a macro has already been defined, whether an optional argument was left empty, or whether two pieces of user input match, all without needing anything beyond TeX's own primitives.

The \texttt{etoolbox} package builds a more approachable layer on top of these primitives: \texttt{\textbackslash{}ifdefined}, \texttt{\textbackslash{}ifboolexpr}, and a family of list- and string-testing conditionals that read closer to their intent than the equivalent primitive expression would. For an author who does not need to reach past this layer, \texttt{etoolbox}'s conditionals are usually the better choice on readability grounds alone; understanding what they expand to in terms of the primitives from the previous section remains valuable for the cases where \texttt{etoolbox}'s conditional does not quite cover what is needed and a primitive-level conditional has to be written by hand.

\hypertarget{loops}{%
\section{Loops}\label{loops}}

TeX has no native looping construct beyond the recursive self-reference a macro can make to its own name, described in Chapter 2 as a direct consequence of the rewriting model rather than a separately implemented feature. \texttt{etoolbox} provides \texttt{\textbackslash{}forcsvlist} and related commands for iterating over a comma-separated list, which covers the overwhelming majority of looping an author actually needs in document preparation: applying the same formatting to each item in a list of keywords, processing each file in a list of includes, or repeating a pattern once per entry in a small, explicitly given set. For anything resembling a numeric loop, incrementing a counter and repeating a body a fixed number of times, \texttt{\textbackslash{}prg\_replicate:nn} from the LaTeX3 kernel, or the older \texttt{\textbackslash{}@whilenum} from the LaTeX2e kernel, provide the same capability at different levels of the toolchain, with the LaTeX3 version generally preferable in new code for its more predictable expansion behavior.

The practical guidance worth taking from this section is restraint: a loop that runs once per chapter to build a table of contents entry, or once per item in a short, fixed list, is a reasonable and common thing to write. A macro-level loop standing in for what would be a straightforward script-level loop over external data, such as generating a hundred entries from a spreadsheet, is usually better handled by generating the LaTeX source for those hundred entries with an external script and including the result, rather than by pushing that iteration into TeX's comparatively awkward looping tools. Knowing where that line sits, between what belongs inside the document's own macros and what belongs in a preprocessing step outside LaTeX entirely, is itself part of a superuser's judgment.

\hypertarget{key-value-interfaces}{%
\section{Key-Value Interfaces}\label{key-value-interfaces}}

A macro with more than two or three optional behaviors quickly becomes unusable if each behavior is controlled by its own positional argument, since the calling code has to remember what each position means and provide placeholder values for anything left at its default. Key-value syntax solves this by letting each option be named explicitly at the call site: \texttt{\textbackslash{}mycommand{[}color=blue,\ bold=true{]}\{text\}} states directly what each setting controls, in any order, with unneeded options simply omitted rather than filled in with a placeholder.

\texttt{pgfkeys}, originally built for PGF and TikZ but usable independently, and \texttt{l3keys}, the LaTeX3 kernel's own key-value system, both provide this pattern as a general tool rather than something tied to graphics. Both let a macro author define a tree of keys, each with a default value, a handler describing what happens when the key is set, and, in \texttt{l3keys}, an optional type check on the value provided. Choosing between them is mostly a matter of context: a project already using TikZ heavily has \texttt{pgfkeys} loaded regardless, while a project building on the LaTeX3 kernel throughout benefits from \texttt{l3keys}'s tighter integration with the rest of that toolchain.

Designing a key-value interface well is closer to designing a small API than to writing an ordinary macro, and the same considerations apply. Every key should have a sensible default, so that a call providing none of the optional keys still produces reasonable output; keys should validate their input where practical, so that a typo in a value produces an immediate, legible error rather than a confusing failure several steps removed from its cause; and keys should compose, meaning that setting one key should not silently change the meaning of another unless that interaction is the documented, intended behavior. A key-value interface that violates any of these is usually harder to use than the positional-argument macro it replaced, since it adds the overhead of remembering key names without delivering the benefit of predictable, independent behavior that was the point of moving to keys in the first place.

\hypertarget{an-example-worth-working-through}{%
\section{An Example Worth Working Through}\label{an-example-worth-working-through}}

Consider a macro for typesetting a labeled example box, the kind that recurs throughout a book like this one, that should support an optional title, an optional numbering scheme independent of the book's other counters, and an optional flag for whether the box should be boxed with a visible border or left as plain indented text. Written with positional arguments, this macro's calling convention would depend on argument order and would need placeholder values at every call site that wanted to set the third option without also setting the first two. Written with \texttt{l3keys}, the same macro exposes \texttt{title}, \texttt{numbered}, and \texttt{boxed} as independently settable keys, each with its own default, callable as \texttt{\textbackslash{}examplebox{[}boxed=false{]}\{...\}} without needing to know or provide anything about the other two options. The interface, not merely the implementation behind it, is what makes the macro pleasant to use across a five-hundred-page book where it might be called differently in dozens of places.

The next chapter turns from what a macro's interface looks like to how its internals should be written so that it behaves correctly and predictably regardless of the surrounding context it is used in, closing out Part III before the book moves on to typography and fonts.
